\subsection{Ближайший методологический предшественник}
В работе Self-collaboration Code Generation via ChatGPT~\cite{dong2023selfcollaboration} уже исследованы ролевые инструкции и максимум взаимодействий в разделах~5.2--5.3. Поэтому удаление ролей или изменение взаимодействий не заявляется новизной настоящего исследования. Таблица~\ref{tab:closest-method} разграничивает реализованные воздействия.

\begin{table}[htbp]\centering\small
\caption{Сопоставление с ближайшим предшественником~\cite{dong2023selfcollaboration}, таблицы~3--4 и приложение~B.1. Раунды взаимодействия и новые вызовы модели являются разными величинами.}\label{tab:closest-method}
\begin{tabular}{p{0.20\textwidth}p{0.31\textwidth}p{0.36\textwidth}}\toprule
Аспект & Self-collaboration & Настоящий эксперимент\\\midrule
Обозначения & Ролевые инструкции и инструкции без ролей & Одинаковые предложения об операциях; меняются только префиксы этапов\\
Взаимодействие & Цикл обратной связи; варьируется максимум взаимодействий & Один или три новых вызова; фиксированная последовательность, полная видимая история\\
Дизайн & Раздельные абляции ролей и взаимодействия & Совместная матрица обозначений и вызовов для каждой отобранной задачи\\
Задачи & HumanEval, MBPP и расширенные тесты & Резервная выборка BigCodeBench; эталонные и неправильные контроли до генерации\\\bottomrule
\end{tabular}\end{table}

Добавленная ценность состоит в репликации и расширении с акцентом на проверяемость: совместный дизайн оценивает взаимодействие обозначений и топологии при одном интерфейсе генерации, а связь результатов штатных тестов с проверенными программами и компоненты токенов делают различия качества и ресурсов проверяемыми. Это новые эмпирические условия и набор доказательств, а не новый принцип ролевого сотрудничества или первая абляция ролей.

Для внешнего сравнения наиболее близок к вопросу об обозначениях метод Self-collaboration вместе с вариантом без ролей. Его реализация должна сохранять обратную связь и правило остановки; переименование нашего фиксированного трёхвызовного MR не воспроизводит этот метод. Отдельный пилот авторской реализации на трёх задачах описан в разделе~\ref{sec:scc-feasibility}; внешнее сравнение с достаточной мощностью и превосходство над SCC остаются непроверенными. В этой схеме D задаёт минимальное реализованное сравнение, а SN и MN служат контролями для выделения эффекта обозначений. INoT* отвечает на отдельный вопрос об алгоритме внутреннего диалога и поэтому не рассматривается как ближайший предшественник абляции ролей.

\subsection{Внутренний диалог и репозиторные методы}
